
\documentclass{article}
\usepackage{colortbl}
\usepackage{makecell}
\usepackage{multirow}
\usepackage{supertabular}

\begin{document}

\newcounter{utterance}

\centering \large Interaction Transcript for game `dond', experiment `semi\_en', episode 18 with reflexion{-}llama{-}3.1{-}8b{-}instruct.
\vspace{24pt}

{ \footnotesize  \setcounter{utterance}{1}
\setlength{\tabcolsep}{0pt}
\begin{supertabular}{c@{$\;$}|p{.15\linewidth}@{}p{.15\linewidth}p{.15\linewidth}p{.15\linewidth}p{.15\linewidth}p{.15\linewidth}}
    \# & $\;$A & \multicolumn{4}{c}{Game Master} & $\;\:$B\\
    \hline

    \theutterance \stepcounter{utterance}  
    & & \multicolumn{4}{p{0.6\linewidth}}{
        \cellcolor[rgb]{0.9,0.9,0.9}{
            \makecell[{{p{\linewidth}}}]{
                \texttt{\tiny{[P1$\langle$GM]}}
                \texttt{You have attempted this same task before and did not succeed. These are the plans you wrote after those attempts.} \\
\texttt{Trial \#1: New plan: } \\
\\ 
\texttt{In my previous attempt, I failed to achieve a mutually beneficial agreement, resulting in a score of zero. Upon reflection, I realize that I focused too much on individual items and their values, rather than considering the overall distribution of items and the potential trade{-}offs. } \\
\\ 
\texttt{Next time, I will take a more holistic approach by considering the entire set of items and the potential value of each combination. I will also pay closer attention to the other player's messages and try to identify their priorities and potential trade{-}offs. } \\
\\ 
\texttt{Specifically, I will start by making a more general proposal that takes into account the overall distribution of items, rather than focusing on individual items. For example, I might propose a 50{-}50 split of the items, or a specific combination of items that I think would be beneficial for both of us. } \\
\\ 
\texttt{I will also try to be more proactive in identifying potential trade{-}offs and offering alternatives that benefit both of us. For example, if the other player is interested in the grapes, I might offer to give them the grapes in exchange for something else that I value more. } \\
\\ 
\texttt{By taking a more holistic approach and being more proactive in identifying potential trade{-}offs, I hope to achieve a mutually beneficial agreement and maximize my score.} \\
\\ 
\\ 
\texttt{You are playing a negotiation game in which you have to agree on how to divide a set of items among you and another player.} \\
\\ 
\texttt{Rules:} \\
\texttt{(a) You and the other player are given a set of items. Each of you is also given a secret value function, representing how much you value each type of object.} \\
\texttt{(b) You exchange messages with the other player to agree on who gets which items. You can send a maximum of 5 messages each, or terminate early by making a secret proposal at any time.} \\
\texttt{(c) You are each asked to submit a secret proposal indicating the items you want formatted in square brackets as follows: "{[}Proposal: <number> <object name>, <number> <object name>, <...>{]}"} \\
\texttt{(d) If your proposals are complementary, i.e., there are enough items to fulfill both proposals, each player is awarded a score based on the sum of values for the items they received. Otherwise, both of you get zero points.} \\
\texttt{(e) Your goal is to maximize the score you receive.} \\
\\ 
\texttt{Let us start.} \\
\\ 
\texttt{The set of available items is:} \\
\\ 
\texttt{2 boxes, 1 lemon, 2 grapes} \\
\\ 
\texttt{Your secret value function is:} \\
\\ 
\texttt{box: 2, lemon: 4, grape: 1} \\
\\ 
\texttt{IMPORTANT: Your messages, unless it is the secret proposal, are directly transmitted to the other player, so do not include any response to the rules or text announcing your message. To make a secret proposal, use the indicated format. Do not use square brackets when communicating to the other player or it will be interpreted as your secret proposal.} \\
            }
        }
    }
    & \\ \\

    \theutterance \stepcounter{utterance}  
    & \multicolumn{4}{p{0.6\linewidth}}{
        \cellcolor[rgb]{0.9,0.9,0.9}{
            \makecell[{{p{\linewidth}}}]{
                \texttt{\tiny{[P1$\rangle$GM]}}
                \texttt{I'd like to start by making a general proposal that takes into account the overall distribution of items. I propose a 50{-}50 split of the items. This way, we can both get a fair share of the items and potentially maximize our scores.} \\
\\ 
\texttt{I'm interested in knowing your thoughts on this proposal and what you value most among the available items.} \\
            }
        }
    }
    & & \\ \\

    \theutterance \stepcounter{utterance}  
    & & & \multicolumn{2}{p{0.3\linewidth}}{
        \cellcolor[rgb]{0.9,0.9,0.9}{
            \makecell[{{p{\linewidth}}}]{
                \texttt{\tiny{[GM$|$GM]}}
                \texttt{continue} \\
            }
        }
    }
    & & \\ \\

    \theutterance \stepcounter{utterance}  
    & & & \multicolumn{2}{p{0.3\linewidth}}{
        \cellcolor[rgb]{0.9,0.9,0.9}{
            \makecell[{{p{\linewidth}}}]{
                \texttt{\tiny{[GM$|$GM]}}
                \texttt{I'd like to start by making a general proposal that takes into account the overall distribution of items. I propose a 50{-}50 split of the items. This way, we can both get a fair share of the items and potentially maximize our scores.} \\
\\ 
\texttt{I'm interested in knowing your thoughts on this proposal and what you value most among the available items.} \\
            }
        }
    }
    & & \\ \\

    \theutterance \stepcounter{utterance}  
    & & \multicolumn{4}{p{0.6\linewidth}}{
        \cellcolor[rgb]{0.9,0.9,0.9}{
            \makecell[{{p{\linewidth}}}]{
                \texttt{\tiny{[GM$\rangle$P2]}}
                \texttt{You have attempted this same task before and did not succeed. These are the plans you wrote after those attempts.} \\
\texttt{Trial \#1: New plan: } \\
\\ 
\texttt{In my last attempt, I failed to achieve a mutually beneficial agreement, resulting in a score of zero. Upon reflection, I realize that I focused too much on individual items and their values, rather than considering the overall distribution of items and the potential trade{-}offs. } \\
\\ 
\texttt{Next time, I will take a more holistic approach by considering the total value of the items and the potential trade{-}offs between them. I will also pay closer attention to the other player's proposals and try to find mutually beneficial agreements. } \\
\\ 
\texttt{Specifically, I will start by making a more general proposal that takes into account the overall distribution of items, rather than focusing on individual items. I will also try to find common ground with the other player by looking for areas of agreement and compromise. } \\
\\ 
\texttt{For example, I might propose a trade that involves multiple items, rather than just one or two. I will also be more willing to listen to the other player's concerns and try to find a solution that works for both of us. } \\
\\ 
\texttt{By taking a more collaborative and holistic approach, I hope to achieve a mutually beneficial agreement and maximize my score.} \\
\\ 
\\ 
\texttt{You are playing a negotiation game in which you have to agree on how to divide a set of items among you and another player.} \\
\\ 
\texttt{Rules:} \\
\texttt{(a) You and the other player are given a set of items. Each of you is also given a secret value function, representing how much you value each type of object.} \\
\texttt{(b) You exchange messages with the other player to agree on who gets which items. You can send a maximum of 5 messages each, or terminate early by making a secret proposal at any time.} \\
\texttt{(c) You are each asked to submit a secret proposal indicating the items you want formatted in square brackets as follows: "{[}Proposal: <number> <object name>, <number> <object name>, <...>{]}"} \\
\texttt{(d) If your proposals are complementary, i.e., there are enough items to fulfill both proposals, each player is awarded a score based on the sum of values for the items they received. Otherwise, both of you get zero points.} \\
\texttt{(e) Your goal is to maximize the score you receive.} \\
\\ 
\texttt{Let us start.} \\
\\ 
\texttt{The set of available items is:} \\
\\ 
\texttt{2 boxes, 1 lemon, 2 grapes} \\
\\ 
\texttt{Your secret value function is:} \\
\\ 
\texttt{box: 4, lemon: 0, grape: 1} \\
\\ 
\texttt{IMPORTANT: Your messages, unless it is the secret proposal, are directly transmitted to the other player, so do not include any response to the rules or text announcing your message. To make a secret proposal, use the indicated format. Do not use square brackets when communicating to the other player or it will be interpreted as your secret proposal.} \\
\\ 
\\ 
\texttt{I'd like to start by making a general proposal that takes into account the overall distribution of items. I propose a 50{-}50 split of the items. This way, we can both get a fair share of the items and potentially maximize our scores.} \\
\\ 
\texttt{I'm interested in knowing your thoughts on this proposal and what you value most among the available items.} \\
            }
        }
    }
    & \\ \\

    \theutterance \stepcounter{utterance}  
    & & & \multicolumn{4}{p{0.6\linewidth}}{
        \cellcolor[rgb]{0.9,0.9,0.9}{
            \makecell[{{p{\linewidth}}}]{
                \texttt{\tiny{[GM$\langle$P2]}}
                \texttt{I think a 50{-}50 split is a good starting point, but I'm not sure if it's the most optimal solution. Considering the items we have, I value the boxes the most, as they seem to be the most valuable. I'm not too concerned about the lemon, but I do like the grapes. } \\
\\ 
\texttt{I'm willing to consider a compromise, but I'd like to know more about your priorities. Do you have a specific item in mind that you'd like to get, or are you open to negotiating the terms of the split?} \\
            }
        }
    }
     \\ \\

    \theutterance \stepcounter{utterance}  
    & & & \multicolumn{2}{p{0.3\linewidth}}{
        \cellcolor[rgb]{0.9,0.9,0.9}{
            \makecell[{{p{\linewidth}}}]{
                \texttt{\tiny{[GM$|$GM]}}
                \texttt{continue} \\
            }
        }
    }
    & & \\ \\

    \theutterance \stepcounter{utterance}  
    & & & \multicolumn{2}{p{0.3\linewidth}}{
        \cellcolor[rgb]{0.9,0.9,0.9}{
            \makecell[{{p{\linewidth}}}]{
                \texttt{\tiny{[GM$|$GM]}}
                \texttt{I think a 50{-}50 split is a good starting point, but I'm not sure if it's the most optimal solution. Considering the items we have, I value the boxes the most, as they seem to be the most valuable. I'm not too concerned about the lemon, but I do like the grapes. } \\
\\ 
\texttt{I'm willing to consider a compromise, but I'd like to know more about your priorities. Do you have a specific item in mind that you'd like to get, or are you open to negotiating the terms of the split?} \\
            }
        }
    }
    & & \\ \\

    \theutterance \stepcounter{utterance}  
    & & \multicolumn{4}{p{0.6\linewidth}}{
        \cellcolor[rgb]{0.9,0.9,0.9}{
            \makecell[{{p{\linewidth}}}]{
                \texttt{\tiny{[P1$\langle$GM]}}
                \texttt{I think a 50{-}50 split is a good starting point, but I'm not sure if it's the most optimal solution. Considering the items we have, I value the boxes the most, as they seem to be the most valuable. I'm not too concerned about the lemon, but I do like the grapes. } \\
\\ 
\texttt{I'm willing to consider a compromise, but I'd like to know more about your priorities. Do you have a specific item in mind that you'd like to get, or are you open to negotiating the terms of the split?} \\
            }
        }
    }
    & \\ \\

    \theutterance \stepcounter{utterance}  
    & \multicolumn{4}{p{0.6\linewidth}}{
        \cellcolor[rgb]{0.9,0.9,0.9}{
            \makecell[{{p{\linewidth}}}]{
                \texttt{\tiny{[P1$\rangle$GM]}}
                \texttt{It sounds like we both value the boxes highly, which is a good starting point for our negotiation. I'm also interested in the grapes, but I don't have a strong attachment to the lemon.} \\
\\ 
\texttt{Considering our shared interest in the boxes, I'd like to propose a slight adjustment to the 50{-}50 split. How about we allocate 1 box to you and 1 box to me, and then split the remaining items (lemon and grapes) in a way that benefits both of us? This way, we can both get a box and then focus on the remaining items.} \\
            }
        }
    }
    & & \\ \\

    \theutterance \stepcounter{utterance}  
    & & & \multicolumn{2}{p{0.3\linewidth}}{
        \cellcolor[rgb]{0.9,0.9,0.9}{
            \makecell[{{p{\linewidth}}}]{
                \texttt{\tiny{[GM$|$GM]}}
                \texttt{continue} \\
            }
        }
    }
    & & \\ \\

    \theutterance \stepcounter{utterance}  
    & & & \multicolumn{2}{p{0.3\linewidth}}{
        \cellcolor[rgb]{0.9,0.9,0.9}{
            \makecell[{{p{\linewidth}}}]{
                \texttt{\tiny{[GM$|$GM]}}
                \texttt{It sounds like we both value the boxes highly, which is a good starting point for our negotiation. I'm also interested in the grapes, but I don't have a strong attachment to the lemon.} \\
\\ 
\texttt{Considering our shared interest in the boxes, I'd like to propose a slight adjustment to the 50{-}50 split. How about we allocate 1 box to you and 1 box to me, and then split the remaining items (lemon and grapes) in a way that benefits both of us? This way, we can both get a box and then focus on the remaining items.} \\
            }
        }
    }
    & & \\ \\

    \theutterance \stepcounter{utterance}  
    & & \multicolumn{4}{p{0.6\linewidth}}{
        \cellcolor[rgb]{0.9,0.9,0.9}{
            \makecell[{{p{\linewidth}}}]{
                \texttt{\tiny{[GM$\rangle$P2]}}
                \texttt{It sounds like we both value the boxes highly, which is a good starting point for our negotiation. I'm also interested in the grapes, but I don't have a strong attachment to the lemon.} \\
\\ 
\texttt{Considering our shared interest in the boxes, I'd like to propose a slight adjustment to the 50{-}50 split. How about we allocate 1 box to you and 1 box to me, and then split the remaining items (lemon and grapes) in a way that benefits both of us? This way, we can both get a box and then focus on the remaining items.} \\
            }
        }
    }
    & \\ \\

    \theutterance \stepcounter{utterance}  
    & & & \multicolumn{4}{p{0.6\linewidth}}{
        \cellcolor[rgb]{0.9,0.9,0.9}{
            \makecell[{{p{\linewidth}}}]{
                \texttt{\tiny{[GM$\langle$P2]}}
                \texttt{I like the direction this is heading. Allocating one box to each of us is a good starting point. } \\
\\ 
\texttt{As for the remaining items, I think it would be beneficial to allocate the lemon to you, since you mentioned you don't have a strong attachment to it. This would leave us with the two grapes, which I think would be a good compromise. I'd be willing to give up one of the grapes in exchange for the box, but I'd like to keep the other grape. } \\
\\ 
\texttt{So, to summarize, I propose that we allocate 1 box to you, 1 box to me, the lemon to you, and 1 grape to me. What are your thoughts on this proposal?} \\
            }
        }
    }
     \\ \\

    \theutterance \stepcounter{utterance}  
    & & & \multicolumn{2}{p{0.3\linewidth}}{
        \cellcolor[rgb]{0.9,0.9,0.9}{
            \makecell[{{p{\linewidth}}}]{
                \texttt{\tiny{[GM$|$GM]}}
                \texttt{continue} \\
            }
        }
    }
    & & \\ \\

    \theutterance \stepcounter{utterance}  
    & & & \multicolumn{2}{p{0.3\linewidth}}{
        \cellcolor[rgb]{0.9,0.9,0.9}{
            \makecell[{{p{\linewidth}}}]{
                \texttt{\tiny{[GM$|$GM]}}
                \texttt{I like the direction this is heading. Allocating one box to each of us is a good starting point. } \\
\\ 
\texttt{As for the remaining items, I think it would be beneficial to allocate the lemon to you, since you mentioned you don't have a strong attachment to it. This would leave us with the two grapes, which I think would be a good compromise. I'd be willing to give up one of the grapes in exchange for the box, but I'd like to keep the other grape. } \\
\\ 
\texttt{So, to summarize, I propose that we allocate 1 box to you, 1 box to me, the lemon to you, and 1 grape to me. What are your thoughts on this proposal?} \\
            }
        }
    }
    & & \\ \\

    \theutterance \stepcounter{utterance}  
    & & \multicolumn{4}{p{0.6\linewidth}}{
        \cellcolor[rgb]{0.9,0.9,0.9}{
            \makecell[{{p{\linewidth}}}]{
                \texttt{\tiny{[P1$\langle$GM]}}
                \texttt{I like the direction this is heading. Allocating one box to each of us is a good starting point. } \\
\\ 
\texttt{As for the remaining items, I think it would be beneficial to allocate the lemon to you, since you mentioned you don't have a strong attachment to it. This would leave us with the two grapes, which I think would be a good compromise. I'd be willing to give up one of the grapes in exchange for the box, but I'd like to keep the other grape. } \\
\\ 
\texttt{So, to summarize, I propose that we allocate 1 box to you, 1 box to me, the lemon to you, and 1 grape to me. What are your thoughts on this proposal?} \\
            }
        }
    }
    & \\ \\

    \theutterance \stepcounter{utterance}  
    & \multicolumn{4}{p{0.6\linewidth}}{
        \cellcolor[rgb]{0.9,0.9,0.9}{
            \makecell[{{p{\linewidth}}}]{
                \texttt{\tiny{[P1$\rangle$GM]}}
                \texttt{I think your proposal is a good compromise, and it seems like we're making progress towards a mutually beneficial agreement.} \\
\\ 
\texttt{However, I'd like to suggest a slight modification to your proposal. Since I value the boxes highly, I'm willing to give up the lemon in exchange for the second grape. This way, I can get both boxes and the second grape, which would be a better outcome for me.} \\
\\ 
\texttt{So, I propose that we allocate 1 box to you, 1 box to me, the second grape to me, and the first grape to you. This way, we both get a box, and I get the second grape, which I value more than the lemon.} \\
            }
        }
    }
    & & \\ \\

    \theutterance \stepcounter{utterance}  
    & & & \multicolumn{2}{p{0.3\linewidth}}{
        \cellcolor[rgb]{0.9,0.9,0.9}{
            \makecell[{{p{\linewidth}}}]{
                \texttt{\tiny{[GM$|$GM]}}
                \texttt{continue} \\
            }
        }
    }
    & & \\ \\

    \theutterance \stepcounter{utterance}  
    & & & \multicolumn{2}{p{0.3\linewidth}}{
        \cellcolor[rgb]{0.9,0.9,0.9}{
            \makecell[{{p{\linewidth}}}]{
                \texttt{\tiny{[GM$|$GM]}}
                \texttt{I think your proposal is a good compromise, and it seems like we're making progress towards a mutually beneficial agreement.} \\
\\ 
\texttt{However, I'd like to suggest a slight modification to your proposal. Since I value the boxes highly, I'm willing to give up the lemon in exchange for the second grape. This way, I can get both boxes and the second grape, which would be a better outcome for me.} \\
\\ 
\texttt{So, I propose that we allocate 1 box to you, 1 box to me, the second grape to me, and the first grape to you. This way, we both get a box, and I get the second grape, which I value more than the lemon.} \\
            }
        }
    }
    & & \\ \\

    \theutterance \stepcounter{utterance}  
    & & \multicolumn{4}{p{0.6\linewidth}}{
        \cellcolor[rgb]{0.9,0.9,0.9}{
            \makecell[{{p{\linewidth}}}]{
                \texttt{\tiny{[GM$\rangle$P2]}}
                \texttt{I think your proposal is a good compromise, and it seems like we're making progress towards a mutually beneficial agreement.} \\
\\ 
\texttt{However, I'd like to suggest a slight modification to your proposal. Since I value the boxes highly, I'm willing to give up the lemon in exchange for the second grape. This way, I can get both boxes and the second grape, which would be a better outcome for me.} \\
\\ 
\texttt{So, I propose that we allocate 1 box to you, 1 box to me, the second grape to me, and the first grape to you. This way, we both get a box, and I get the second grape, which I value more than the lemon.} \\
            }
        }
    }
    & \\ \\

    \theutterance \stepcounter{utterance}  
    & & & \multicolumn{4}{p{0.6\linewidth}}{
        \cellcolor[rgb]{0.9,0.9,0.9}{
            \makecell[{{p{\linewidth}}}]{
                \texttt{\tiny{[GM$\langle$P2]}}
                \texttt{I think your proposal is a good one, and it's a fair trade considering our values. I'm willing to give up the first grape in exchange for the lemon, which I don't value as much. } \\
\\ 
\texttt{So, to confirm, we have a deal: 1 box to you, 1 box to me, the second grape to you, and the first grape to me, with the lemon going to you. I think this is a mutually beneficial agreement, and I'm happy to accept it. } \\
\\ 
\texttt{{[}Proposal: 1 box, 1 box, 1 lemon, 1 grape{]}} \\
            }
        }
    }
     \\ \\

    \theutterance \stepcounter{utterance}  
    & & & \multicolumn{2}{p{0.3\linewidth}}{
        \cellcolor[rgb]{0.9,0.9,0.9}{
            \makecell[{{p{\linewidth}}}]{
                \texttt{\tiny{[GM$|$GM]}}
                \texttt{proposal must include every item type only once} \\
            }
        }
    }
    & & \\ \\

    \theutterance \stepcounter{utterance}  
    & \multicolumn{4}{p{0.6\linewidth}}{
        \cellcolor[rgb]{0.95,0.95,0.95}{
            \makecell[{{p{\linewidth}}}]{
                \texttt{\tiny{[P2$\rightarrow$P1]}}
                \texttt{Post{-}episode Reflexion (not sent to GM):} \\
\texttt{New plan: } \\
\\ 
\texttt{In my previous attempt, I failed to achieve a mutually beneficial agreement, resulting in a score of zero. Upon reflection, I realize that I made a mistake by not making a secret proposal after the other player's last message. } \\
\\ 
\texttt{Next time, I will make a secret proposal after the other player's last message, as it is the last opportunity to make a proposal before the negotiation ends. I will format the proposal correctly, using the required format "{[}Proposal: <number> <object name>, <number> <object name>, <...>{]}" and including all the items that I want. } \\
\\ 
\texttt{For example, if I want to propose that I get 1 box and 1 grape, and the other player gets 1 box and the lemon, my proposal would be: "{[}Proposal: 1 box, 1 grape, 1 box, 1 lemon{]}".} \\
            }
        }
    }
    & & \\ \\

    \theutterance \stepcounter{utterance}  
    & & & \multicolumn{4}{p{0.6\linewidth}}{
        \cellcolor[rgb]{0.95,0.95,0.95}{
            \makecell[{{p{\linewidth}}}]{
                \texttt{\tiny{[P1$\rightarrow$P2]}}
                \texttt{Post{-}episode Reflexion (not sent to GM):} \\
\texttt{New plan: } \\
\\ 
\texttt{In my last attempt, I failed to achieve a mutually beneficial agreement, resulting in a score of zero. Upon reflection, I realize that I made a secret proposal without following the correct format. } \\
\\ 
\texttt{Next time, I will make sure to follow the correct format for a secret proposal, which is to use square brackets and the specified format: "{[}Proposal: <number> <object name>, <number> <object name>, <...>{]}"} \\
            }
        }
    }
     \\ \\

\end{supertabular}
}

\end{document}
